\documentclass[12pt]{article}
\usepackage{amssymb}
\usepackage{amsmath}
\usepackage{color}
\usepackage{graphicx}
\usepackage{mdframed}
\usepackage{listings, xcolor}
\usepackage{booktabs}
% \usepackage[dvipsnames]{xcolor} % Non va!
\usepackage{textcomp}
\usepackage[table]{xcolor}
\newcommand{\gr}{\rowcolor[HTML]{EFEFEF}}
\lstdefinelanguage{myC}{
	keywords=[1]{break, case, continue, default, do
		, else, false, for, if, const, return, switch, true, while}, % generic keywords
	keywordstyle=[1]\color{blue}\bfseries,
	keywords=[2]{bool, int, long, float, double, byte, short, char, void, signed, unsigned}, % types
	keywordstyle=[2]\color{teal}\bfseries,
	keywordstyle=[2]\color{violet}\bfseries,
	keywords=[3]{NULL},
	keywordstyle=[3]\color{violet}\bfseries,
	identifierstyle=\color{black},
	sensitive=false,
	comment=[l]{//},
	morecomment=[s]{/*}{*/},
	commentstyle=\color{violet}\ttfamily,
	commentstyle=\normalsize\ttfamily\color{grey}, % scriptsize
	stringstyle=\color{forestgreen}\ttfamily,
	morestring=[b]',
	morestring=[b]"
}

\lstset{
	language=myC,
	backgroundcolor=\color{veryLightgray},
	extendedchars=true,
	basicstyle=\small\ttfamily,
	showstringspaces=false,
	showspaces=false,
	numbers=none,
	numberstyle=\normalsize,
	numbersep=9pt,
	tabsize=2,
	upquote=true,
	breaklines=true,
	showtabs=false,
	captionpos=b
}

\definecolor{myBlue}{rgb}{0.5,0.5,1}
\definecolor{myLightBlue}{rgb}{0.35,0.6,0.8}
\definecolor{myBlack}{rgb}{0,0,0}
\definecolor{myGreen}{rgb}{0.1,0.6,0.2}
\definecolor{myGray}{rgb}{0.5,0.5,0.5}
\definecolor{myLightgray}{rgb}{0.95,0.95,0.95}
\definecolor{verylightgray}{rgb}{.97,.97,.97}
\definecolor{myMauve}{rgb}{0.58,0,0.82}
\definecolor{forestgreen}{rgb}{0.13, 0.55, 0.13}
\definecolor{forestgreen}{rgb}{0.13, 0.55, 0.13}
\lstdefinelanguage{customc}{
    language=C,
    backgroundcolor = \color{myLightgray},
    basicstyle=\small\ttfamily\color{myBlack},
    keywordstyle=\color{myLightBlue},
    keywordstyle=[2]\color{red},
    commentstyle=\small\ttfamily\color{myGreen},
    morekeywords={RequirePackage,ProvidesPackage},
    %
    % The special highlighting works for '!if', '!endif' and '!else'
    % But it doesn't work for '#if', '#endif' and '#else'.
    alsoletter = {!},
    keywords=[2]{!if,!endif,!else},
    otherkeywords={define,include,\# }
}

\lstdefinestyle{myCustomc}{
    language = customc,
    % keywordstyle = \color{myMauve},
}

\lstset{escapechar=@,style=myCustomc}


\definecolor{grey}{rgb}{0.3,0.3,0.3}
\definecolor{lightgrey}{rgb}{0.9,0.9,0.9}

\thispagestyle{empty}
\setlength{\textwidth}{18.5cm}
\setlength{\topmargin}{-2.5cm}
\setlength{\textheight}{24.5cm}
\setlength{\oddsidemargin}{-1cm}
\setlength{\evensidemargin}{-1cm}

\begin{document}
\begin{center}
  {\LARGE Compito di Programmazione - Bioinformatica}\\
  {9 settembre 2025 (tempo disponibile: 2 ore)}\\
  {(si consegnino i file \texttt{sentence.h} e \texttt{sentence.c} opportunamente completati)}
\end{center}

La struttura \texttt{sentence} implementa una frase, cio\`e una lista di parole. \`E definita
dentro \texttt{sentence.h}:

\begin{center}
  \begin{lstlisting}[language=myC]
    struct sentence {
      char *word;
      struct sentence *next;
    };
  \end{lstlisting}
\end{center}

Una frase pu\`o essere costruita a partire da un array di stringhe
(\texttt{construct\_sentence})
e successivamente deallocata (\texttt{free\_sentence}).
Inoltre se ne pu\`o calcolare la lunghezza, cio\`e il numero di caratteri
totale delle sue parole, contando anche uno spazio tra una parola
e l'altra (\texttt{length\_sentence}). Infine, una
frase pu\`o essere trasformata in una nuova stringa,
ottenuta concatenando le sue parole e aggiungendo uno spazio tra una parola
e l'altra (\texttt{concat\_sentence}). La stessa operazione pu\`o essere
effettuata in senso inverso, cio\`e iniziando, nella concatenazione,
dall'ultima parola e terminando con la prima parola (\texttt{inverse\_concat\_sentence}).

Il seguente \texttt{main.c}, gi\`a fatto, completo e da non modificare, esemplifica
l'utilizzo di una frase e delle sue funzioni citate sopra:

\begin{mdframed}[backgroundcolor=verylightgray]
  \begin{lstlisting}[language=myC]
#include <stdio.h>
#include <stdlib.h>
#include "sentence.h"

int main(void) {
  char *words[] = {
    "Nel", "mezzo", "del", "cammin", "di", "nostra", "vita"
  };

  struct sentence *s = construct_sentence(words, 7);
  printf("La lunghezza della frase e' %i\n", length_sentence(s));
  char *concat = concat_sentence(s);
  printf("La sua concatenazione e' \"%s\"\n", concat);
  free(concat);
  concat = inverse_concat_sentence(s);
  printf("La sua concatenazione inversa e' \"%s\"\n", concat);
  free(concat);
  free_sentence(s);

  return 0;
}
  \end{lstlisting}
\end{mdframed}

\noindent
Quando avrete implementato tutte le funzioni correttamente,
l'esecuzione del \texttt{main} dovr\`a stampare:

\begin{mdframed}[backgroundcolor=verylightgray]
\begin{verbatim}
La lunghezza della frase e' 35
La sua concatenazione e' "Nel mezzo del cammin di nostra vita"
La sua concatenazione inversa e' "vita nostra di cammin del mezzo Nel"
\end{verbatim}
\end{mdframed}

\vspace*{1ex}
\begin{center}{\Large Esercizio 1} ($1$ punto)\end{center}
Si completi il file \texttt{sentence.h} aggiungendo la dichiarazione delle cinque funzioni
\texttt{construct\_sentence}, \texttt{free\_sentence},
\texttt{length\_sentence}, \texttt{concat\_sentence} e
\texttt{inverse\_concat\_sentence}. Dentro \texttt{sentence.c} trovate
la bozza di tali funzioni.

\vspace*{1ex}
\begin{center}{\Large Esercizio 2} ($4$ punti)\end{center}
Si completi la funzione \texttt{construct\_sentence} dentro \texttt{sentence.c}.
Si noti che tale funzione copia, dentro la
frase risultante, il puntatore alle stringhe passate come argomento, quindi
non fa una copia di tali stringhe.

\vspace*{1ex}
\begin{center}{\Large Esercizio 3} ($3$ punti)\end{center}
Si completi la funzione \texttt{free\_sentence} dentro \texttt{sentence.c}.
Si noti che tale funzione deve deallocare tutta la lista, precedentemente allocata
con \texttt{construct\_sentence}, non solo il suo primo elemento.

\vspace*{1ex}
\begin{center}{\Large Esercizio 4} ($5$ punti)\end{center}
Si completi la funzione \texttt{length\_sentence} dentro \texttt{sentence.c}.
Si noti che tale funzione, nel calcolare la lunghezza,
considera uno spazio tra una parola e l'altra delle frase.

\vspace*{1ex}
\begin{center}{\Large Esercizio 5} ($8$ punti)\end{center}
Si completi la funzione ausiliaria \texttt{concat\_sentence\_rec} dentro \texttt{sentence.c},
che concatena le parole di una frase dentro una stringa risultato, con uno spazio tra una
parola e l'altra. \textbf{Tale funzione deve essere ricorsiva}. La si usi quindi per
completare la funzione \texttt{concat\_sentence} dentro \texttt{sentence.c}.
\textbf{Suggerimento:} in \texttt{concat\_sentence\_rec} si considerino i casi:
frase vuota, frase di una sola parola, frase di pi\`u parole.

\vspace*{1ex}
\begin{center}{\Large Esercizio 6} ($10$ punti)\end{center}
Si completi la funzione ausiliaria \texttt{inverse\_concat\_sentence\_rec} dentro \texttt{sentence.c},
che concatena le parole di una frase dentro una stringa risultato, con uno spazio tra una
parola e l'altra, in ordine inverso, cio\`e mettendo all'inizio l'ultima parola e alla fine la prima parola.
\textbf{Tale funzione deve essere ricorsiva}. La si usi quindi per
completare la funzione \texttt{inverse\_concat\_sentence} dentro \texttt{sentence.c}.
\textbf{Suggerimento:} in \texttt{inverse\_concat\_sentence\_rec}
si considerino i casi: frase vuota, frase di una sola parola, frase di pi\`u parole.

\end{document}
